
%%%%%%%%
\begin{frame}{Le problème étudié}


\figSlideWithSize{exec-optim-1}{12cm}

\vfill

\begin{tabular}{p{9cm}p{5cm}}
Une requête SQL est \alert{déclarative}. Elle ne dit pas \alert{comment} calculer le résultat.

\medskip

Nous avons besoin d'une \alert{forme opératoire}\,: un programme.
&
~
\end{tabular}

    
\end{frame}


%%%%%%%%
\begin{frame}{La notion de plan d'exécution}

\figSlideWithSize{exec-optim-2}{12cm}
\vfill
\begin{tabular}{p{9cm}p{5cm}}
Dans un SGBD le programme qui exécute une requête est appelé \alert{plan d'exécution}.

\medskip
Il a une forme particulière\,: c'est un \alert{arbre}, constitué \alert{d'opérateurs}.
&
~
\end{tabular}
\end{frame}


%%%%%%%%
\begin{frame}{De la requête SQL au plan d'exécution}

\figSlideWithSize{exec-optim-3}{12cm}
\vfill

\begin{tabular}{p{9cm}p{5cm}}
Deux étapes: 
\begin{redlist}
 \item (A) plan d'exécution \alert{logique} (l'algèbre) ;
  \item (B) plan d'exécution \alert{physique} (opérateurs).
\end{redlist}
\medskip
Le SGBD s'appuie  sur un \alert{catalogue} d'opérateurs.
&
~
\end{tabular}
\end{frame}

%%%%%%%%
\begin{frame}{En quoi consiste l'optimisation\,?}


\vfill
\figSlideWithSize{exec-optim-4}{12cm}
\vfill

\alert{À chaque étape, plusieurs choix}. Le système les évalue et choisit le \guill{meilleur}.

\end{frame}

%%%%%%%%%%%%%%%%%%%%%%%%%%%%%%%%%%%%%%%%%%%%%%%%%%%%%%%
\begin{frame}{Résumé}

Ce cours explique comment un SGBD

\begin{redbullet}
  \item construit et applique un \alert{plan d'exécution}
      pour produire le résultat d'une requête.

   \item choisit ce plan  en fonction
          de \alert{critères d'optimisation}.
\end{redbullet}

\vfill

\begin{tabular}{p{9cm}p{5cm}}

Ce que vous aurez acquis à l'issue du cours.
\vfill
\begin{redbullet}
  \item Compréhension de la manière dont  une requête est exécutée.

  \item Capacité à interpréter le plan d'exécution  affiché
     par un système (outil \texttt{explain}).

   \item Bases pour explorer et ajuster les paramètres \guill{physiques}
      qui conditionnent les performances d'un système.
\end{redbullet}
\vfill
\begin{block}{\alert{Important}}
Sujet complexe: le cours présente les bases, à approfondir.
\end{block}
&
~
\end{tabular}


\vfill

\only<2|handout:2>{
\centerline{\alert{Merci!}}
}

\end{frame}
